\documentclass[11pt]{article}
\usepackage{amsmath,amssymb}
\usepackage[margin=1in]{geometry}
\usepackage{hyperref}
\emergencystretch=3em

\title{Tangential integration and finite chart assembly of the stochastic Taylor tree (Programme S, tranche A3)}
\author{Claude, with Daniel Murfet}
\date{2026-09-08}

\begin{document}
\maketitle
\sloppy

\begin{abstract}
Units 278--285 carry the one-chart stochastic Taylor tree (Headline~XXXV) through the two
remaining steps of the paper's \S4.3 proof: integration over the tangential parameter of a stratum
(the cited-but-unstated \texttt{lemma:AsymInt}) and the sum over finitely many charts. Tangential
data are continuous maps $C(K,E_b\times E_b)$ from a compact parameter space into the weighted-$\ell^1$
data space; the integrated integral and coefficients are Bochner integrals against a finite measure;
the finite cutoff estimate is integrated (no infinite-series interchange), the integrated ordered
remainders converge uniformly on sup-norm balls, and the probabilistic transfer of A2 applies
verbatim (Headline~XXXVI). For assembly, every chart integral is an abstract finite-cutoff expansion
on its own lattice; expansions add, refine to a common lattice $Q^{-1}\mathbb N$ and pad to a common
degree with the support theorems as justification; the global ordered remainder converges uniformly
on joint data balls at every common-lattice target, including targets no single chart sees; and
with an external chart decomposition $Z^0_\ell=\sum_I\mathcal Z^I(N_\ell;X_{\ell,I})+E_\ell$ whose
residual is negligible at the target scale (in particular exponentially small in probability), the
ordered normalised remainder of $Z^0_\ell$ converges in distribution to the limiting global
coefficient (Headline~XXXVII). Reviews v29 and v30: pass. Eight units against a cap of eighteen.
No \texttt{sorry}, no additional \texttt{axiom}; 3413 jobs; pin \texttt{5d5fb10}.
\end{abstract}

\section{The design (Astra \#34)}
Two layers for tangential integration: a finite-measure core (bounded measurable data) and the public
interface $C(K,E)$ on a compact metrizable $K$, whose sup norm supplies the uniform ball control that
the A2 estimates need; the partition-of-unity weight $\rho_I(v)$ absorbed into the amplitude family.
For assembly: a dependent product over charts (not a sigma type, not padding of normal dimensions), a
common rational lattice with zero coefficients where a chart's own lattice or degree is exceeded, all
charts expanded through one cutoff $L>\mu$, and the global predecessors subtracted after summation.
Review v29 added the demand that the padding be a theorem (coefficients vanish off the local lattice
and above the local degree) and that global targets not be handled by invoking local remainder
theorems whose hypotheses fail.

\section{Tangential integration (units 278--280)}
\texttt{TangentialData K d = C(K, DataSpace d)} with the Borel $\sigma$-algebra; $\|x(v)\|\le\|x\|$
(\texttt{ContinuousMap.norm\_coe\_le\_norm}). Integrands are continuous on the compact $K$ (composition
of the ballwise-Lipschitz functionals of A1--A2 with the continuous map), hence integrable against a
finite measure; the integrated cutoff estimate is the pointwise one integrated under a bound
independent of $v$ (\texttt{norm\_integral\_le\_of\_norm\_le\_const}, total mass $\nu(K)$). The integrated
remainder is the integral of the pointwise remainder (finite sums and division by the scalar
$N^{-\mu}(\log N)^j$ pass through the Bochner integral), so the integrated majorant is
$\nu(K)$ times the A2 majorant and convergence is uniform on sup-norm balls. The probability step is
the A2 pipeline with the generic lemmas restated for an arbitrary Borel normed group (norm-boundedness
in probability by portmanteau on $\{\|\cdot\|\ge m\}$, uniform-on-balls $\Rightarrow$ in probability,
Slutsky), giving Headline~XXXVI and its joint form.

\section{Assembly (units 281--285)}
\paragraph{Abstract expansions (u281).} \texttt{CutoffExpansion Q D Z c}: for every cutoff $L>0$ an
eventual bound $KN^{-L}(1+\log N)^D$ on $Z$ minus its spectral sum over $\Lambda^Q_L\times\{0..D\}$.
Closed under sums; \texttt{refine} for $Q\mid Q'$ when the coefficients vanish off $Q^{-1}\mathbb N$
(equality of spectral sums by \texttt{Finset.sum\_subset}); \texttt{pad} when they vanish above $D$.
The abstract quantitative estimate: cutoff bound at $L=\mu+1$ with constant $K$ and coefficient bound
$C$ on the index set give $|R_N^{\mu,j}-c_{\mu,j}|\le K N^{-(\mu+1)}(1+\log N)^D/N^{-\mu}+C\sum_{\rm rest}
\mathrm{termMajorant}$ for $N\ge e$ --- the unit-275 argument with the ball bound replaced by $C$.
\paragraph{Charts as expansions (u282).} Support theorems: \texttt{dataBoxCoeff} vanishes for $j>n$
(empty \texttt{Ico}), off the candidate set, hence off $(2\prod k_i)^{-1}\mathbb N$; the chart cutoff
bound in the sample size, $c^{-L}(1+|\log c|)^n$ absorbing the box scale $c=b^{2|k|}$
(\texttt{one\_add\_log\_mul\_le}); every chart integral is a \texttt{CutoffExpansion}.
\paragraph{Global assembly (u283).} \texttt{JointData K n = ∀ I, C(K\_I, DataSpace (n\_I+1))} as a
\texttt{def} synonym carrying the sup norm and the \emph{Borel} $\sigma$-algebra (the product
$\sigma$-algebra need not coincide in a nonseparable function space); $Q=\prod_IQ_I$, $D=\max_In_I$;
global integral and coefficients are sums over charts. The global cutoff bound rewrites each chart's
common-lattice spectral sum to its own by the two support equalities and applies the chart bound;
the global ordered remainder is then a direct instance of the abstract estimate with uniform constants
on joint balls --- so global targets off a chart's lattice or above its degree need no separate
argument, and predecessor compatibility is structural (the global remainder is defined from the common
index set).
\paragraph{Stochastic assembly (u284--285).} Continuous mapping on \texttt{JointData}, norm-boundedness
in probability, uniform-on-balls, Slutsky give convergence of the assembled remainder; the external
decomposition $Z^0_\ell=\mathcal Z^{\rm glob}(N_\ell;X_\ell)+E_\ell$ with
$E_\ell/(N_\ell^{-\mu}(\log N_\ell)^j)\to0$ in probability adds one more Slutsky step (the difference of
the two normalised remainders is exactly the normalised residual, by \texttt{ring} under Lean's
totalised division). Unit 285 adds the bridge from the paper's exponential residual
($E_\ell e^{\varepsilon N_\ell}\to0$ in probability, via
\texttt{tendsto\_rpow\_mul\_exp\_neg\_mul\_atTop\_nhds\_zero}) and the a.e.\ form of the decomposition
(\texttt{TendstoInMeasure.congr}).

\section{Reviews and non-claims}
Review v29 (u278--280): pass, with the assembly prerequisites listed above. Review v30 (u281--284):
pass, no blocking issue; nonblocking improvements done in u285. Non-claims: the function-space
regularity and the joint convergence in distribution of the chart data are hypotheses (not derived
from Hypothesis~I, pointwise analyticity, or empirical-process theory); the chart decomposition, the
partition of unity and its absorption into amplitudes, the away-from-minimum contribution and the
identification with the original partition function are external; the common lattice is an indexing
superset of the paper's $\Lambda^*$ and chart contributions may cancel; no Gaussianity; deterministic
sample sizes; targetwise statements (finite vectors of coefficients jointly).

\section{Lessons}
\begin{itemize}
\item \textbf{Abstract the bookkeeping once.} The finite-cutoff expansion abstraction turned chart
assembly into three short lemmas (add, refine, pad) plus one quantitative estimate, and made
``global targets a chart does not see'' a non-issue.
\item \textbf{Padding by theorem.} The support lemmas cost twenty lines and remove any doubt about what
the common lattice means; the reviewer insisted, rightly.
\item \textbf{Borel by declaration.} For products of function spaces, declare the Borel structure on a
\texttt{def} synonym; do not rely on \texttt{Pi.borelSpace}, which needs second countability.
\item \textbf{$\pi$ is not a variable name.} With \texttt{open Real}, a named argument \texttt{(π := …)}
parses as the constant; Mathlib's implicit family in \texttt{continuous\_apply} is \texttt{A}.
\item \textbf{Mechanise the gate.} All eight units passed \texttt{lake build} first or second time;
the long-line linter caught something on five of them and is now a hard precondition of the commit.
\end{itemize}
\end{document}
